% !TeX root = ../main.tex
% !TeX program = xelatex
% !TeX encoding = UTF-8
\section{研究问题与技术贡献}

供热转型同时涉及能源替代、既有设施、用热舒适和电网接入。工业余热受到温位和输配约束，分散取暖取决于低温设备能力，夏热冬冷地区需要统筹冷热服务。规划因而必须同时回答资源如何分配、建筑何时用热，以及调度能否满足设备和网络边界。

依据比赛对模型、逐时仿真、多目标、多尺度及代码的要求\supcite{1}，本文以乌鲁木齐、固原和武汉为背景，围绕“物理一致的方案生成—多目标选择—逐时及工程校核”组织研究，并以公共数据对照实验检验方法。

本文的特色集中在四个方面。第一，将共享余热、独立储热和建筑热惯性接入同一能量与状态约束体系，显式限制同一资源的重复分配，并通过区域公用电上限连接建筑调度。第二，将实测负荷的分布保真与能量守恒共同用于代表日权重校正，配合峰值、爬坡和守恒门槛判断压缩方案能否采用。第三，在固定设备研究中使用厂家制热量与输入功率的有效性能域，在网络研究中以校核失败模式的排除推动重新优化，使方案选择响应设备和输配边界。第四，将费用、能量、碳排与求解界分别留痕，给出可核查的条件收益及其代价。

上述贡献体现为模型集成、约束实现与对照实验：以既有热网络、RC及规划方法为基础，通过统一接口连接，并分别验证其作用。

\section{数据基础与三城实验设计}
\subsection{数据来源和用途}

数据按气象驱动、建筑原型、实测计量、设备性能和费用因子分层组织。各层通过对象标识、时间戳与单位映射连接，保留原始文件及变换记录，使模型参数与结果可逐项追溯。表\ref{tab:data-scope}给出数据分工。

\begin{table}[htbp]
\centering
\caption{研究数据及其用途}\label{tab:data-scope}
\begin{tabularx}{\textwidth}{@{}>{\raggedright\arraybackslash}p{0.23\textwidth}X >{\raggedright\arraybackslash}p{0.23\textwidth}@{}}
\toprule
数据依据 & 在本文中的作用 & 证据性质\\
\midrule
NASA POWER逐时气象\supcite{9} & 三城相同窗口下的气候输入；登记时钟、单位及转换 & 公开格点产品\\
BESTEST 600原型\supcite{11} & 建筑几何和热工基准的来源；按登记方法派生1R1C参数 & 公共原型的模型变换\\
BDG2建筑计量\supcite{10} & 六栋建筑的负荷压缩与代表日权重对照 & 外部计量基准\\
TICA产品样本\supcite{12} & 两型号共165个制热量/功率点及缺测角点 & 厂家公开工况表\\
电价及排放因子\supcite{13,14} & 统一比较费用和参考碳排；保留适用条件 & 共同条件性核算\\
OSM空间数据 & 三城建筑与道路候选的空间参考 & 公开空间数据\\
\bottomrule
\end{tabularx}
\end{table}

三城气候对照采用北京时间2024年11月16日00:00至2025年3月16日00:00的左闭右开窗口，共120日、2880个连续小时。建筑与服务条件保持一致，以识别气象影响。费用比较统一采用湖北条件性工商业电价及固定平均排放因子，碳排为同一平均因子下的参考核算。厂家热泵算例将48~m$^2$公共原型合成放大20倍至960~m$^2$，温度边界为20--27~$^\circ$C，初末状态均为20~$^\circ$C。

本文采用两层验证：三城公共原型实验检验同条件调度机制，外部计量实验检验负荷压缩质量。所用连续窗口、资产与边界均随实验登记。

\subsection{三类场景与验证分工}

乌鲁木齐对应既有集中供热与工业余热接入问题，重点是资源共享、温位适配和远端输配；固原对应清洁取暖替代，重点是设备低温能力与接入约束；武汉对应冷热需求与新建系统，重点是建筑热惯性和区域协同。三城均参与统一气候窗口和设备调度对照；网络反馈采用登记的72小时参考系统，公用电耦合采用武汉气象下的双建筑算例。各组实验分别识别气候、设备、建筑柔性和输配约束的作用。

负荷压缩采用BDG2境外建筑计量作为外部基准。三城OSM空间数据承担建筑定位与候选路径表达，资源池通过对象标识关联热源能力和服务对象，实现空间信息与能量约束的分工。

\subsection{统一比较与计量规则}

同一对照固定天气、资产、服务温度、初末状态和计费条件，仅改变所研究的参数或约束。功率、能量、费用与排放分别用kW、kWh、元与kg计量。代表日权重用于汇总，连续状态在原时序中检验；时间轴检查器审核年份、时区和小时连续性。每个结果保存配置、输入指纹、求解状态及独立账本，保证比较可复核。
